\section{Agent Eval Taxonomy：能力、应用与广覆盖评估}

基础任务类型回答“如何评分”，taxonomy 回答“评分覆盖什么”。课程将 Agent eval 分成三层：specific capability、specific application、general set of applications。三层不是高低等级，而是不同诊断粒度。能力评估便于定位工具调用、规划或记忆缺口；应用评估检验完整工作流；广覆盖套件用于比较通用 Agent 的能力轮廓。

\subsection{Specific capability：把复杂 Agent 拆成可诊断能力}

slides 27--30 汇总 survey 中的能力类别。规划/推理、function calling/tool use、self-reflection、LLM-based evaluation、memory 等都可单独构造任务。能力 benchmark 的优势是控制变量清楚：例如 BFCL 关注函数选择和参数，memory benchmark 关注跨轮信息保留。缺点是模块分数不能自动预测端到端成功，真实应用会把多个能力以时序方式耦合。

\lecturefigure{slides-images/slide-027.png}{官方 slide 27：规划与推理类 Agent capability benchmark}
\lecturefigure{slides-images/slide-028.png}{官方 slide 28：Function Calling 与 Tool Use benchmark}
\lecturefigure{slides-images/slide-029.png}{官方 slide 29：Self-Reflection 与 LLM-based evaluation capability}
\lecturefigure{slides-images/slide-030.png}{官方 slide 30：Memory capability benchmark}

\begin{knowledgebox}{读图：能力名称背后要落到可观察行为}
“规划”可以测子目标分解、顺序约束或重规划；“工具使用”可测工具选择、参数、时机和错误恢复；“记忆”可测写入、检索、更新与遗忘。若 benchmark 只用任务名称分类，却没有明确失败判据，不同论文的同名分数不可比较。能力 eval 应输出细分错误，而不只是一个 aggregate。
\end{knowledgebox}

\begin{warningbox}{模块高分不等于组合可靠}
Agent 可能分别通过 tool-use 与 planning benchmark，却在真实任务中因计划与工具状态不同步失败。组合系统还会出现错误传播、预算竞争和权限冲突。能力 benchmark 适合诊断和回归，应用 benchmark 才能检查接口契约和长链路可靠性。
\end{warningbox}

\subsection{Specific application：用真实工作流检验系统闭环}

slides 32--36 覆盖 web agents、software engineering、scientific agents、conversational agents 以及医疗、教育、法律、管理和金融等应用。应用评估的任务、工具、数据和用户目标更接近真实环境，能够暴露流程错误；但每个 domain 有自己的安全、许可和专业标准，分数可迁移性有限。

\lecturefigure{slides-images/slide-032.png}{官方 slide 32：Web Agent benchmark}
\lecturefigure{slides-images/slide-033.png}{官方 slide 33：Software Engineering Agent benchmark}
\lecturefigure{slides-images/slide-034.png}{官方 slide 34：Scientific Agent benchmark}
\lecturefigure{slides-images/slide-035.png}{官方 slide 35：Conversational Agent benchmark}
\lecturefigure{slides-images/slide-036.png}{官方 slide 36：医疗、教育、法律、管理与金融应用}

\begin{knowledgebox}{应用评估的四层真实性}
任务真实性：目标是否来自真实需求；环境真实性：工具、数据和约束是否接近生产；行为真实性：模型是否需要完整交互而非一次生成；后果真实性：错误成本和安全边界是否被表达。为了可复现，benchmark 常简化环境，因此“realistic”应具体说明保留和省略了什么。
\end{knowledgebox}

\textbf{课堂提示：}应用 benchmark 容易被误用为行业准入证明。课程列出大量论文，是为了展示评估空间，而不是说在某个医疗或法律 benchmark 得分高就可以部署。高风险领域还需专家审查、真实流程验证、监管合规和持续监控。

\subsection{General set：广覆盖不是简单求平均}

能力和应用评估各自聚焦局部，general set 则试图回答模型跨任务迁移是否稳定。本小节关注任务权重和汇总方法：如果没有能力剖面、难度分桶与成本信息，一个总分会把“少数领域很强、关键领域失效”的系统包装成通用 Agent。
slide 38 展示 general Agent benchmark 试图跨任务/应用测量通用能力。这类套件适合比较 Agent 系统的能力轮廓，但任务权重会隐式定义什么叫“通用”。如果 coding 占比过高，综合分会偏向代码模型；若简单任务很多，平均值会掩盖难题失败。

\lecturefigure{slides-images/slide-038.png}{官方 slide 38：General Agent Evaluation 的跨应用集合}

\begin{importantbox}{综合分必须带能力剖面}
报告总体分时同时给出 domain、难度、工具、轨迹长度和成本分桶；说明权重与缺失数据处理；避免用平均值掩盖安全关键失败。通用 Agent 的目标不是每项相同，而是在任务分布变化时保持可预测、可诊断的性能。
\end{importantbox}

\subsection{本章小结}

能力层负责定位，应用层负责闭环，general set 负责广覆盖比较。一个成熟 eval program 通常三层并存：快速 capability regression、少量高保真 application scenario，以及周期性的 broad suite。

